% ============================================================================
% Public CV -- AI engineering (industry)
%
% This is the PUBLIC version of the CV that is downloadable from
% ben-whewell.com/resume.html. It deliberately contains NO phone number and NO
% email address; readers are pointed to the request link on the website
% instead. Do not add direct contact details to this file.
%
% Build (TinyTeX / TeX Live):
%   pdflatex -interaction=nonstopmode -jobname=Ben_Whewell_CV cv_ai_engineering.tex
% ============================================================================
\documentclass[10.5pt,letterpaper]{article}

\usepackage{lmodern}
\usepackage[T1]{fontenc}
\usepackage[utf8]{inputenc}
\usepackage[margin=0.62in,bottom=0.55in,top=0.55in]{geometry}
\usepackage{microtype}
\usepackage{enumitem}
\usepackage{titlesec}
\usepackage{xcolor}
\usepackage[hidelinks]{hyperref}

\definecolor{accent}{RGB}{5,68,104}
\definecolor{muted}{RGB}{95,95,95}

\hypersetup{colorlinks=true,urlcolor=accent,pdftitle={Benjamin Whewell -- CV},
  pdfauthor={Benjamin Whewell},pdfsubject={AI/ML engineering CV}}

\pagestyle{empty}
\setlength{\parindent}{0pt}
\setlength{\parskip}{0pt}

% --- section headings -------------------------------------------------------
\titleformat{\section}{\normalfont\large\bfseries\color{accent}}{}{0pt}{\MakeUppercase}[{\vspace{2pt}\hrule height 0.6pt}]
\titlespacing{\section}{0pt}{7pt}{4pt}

% --- helpers ---------------------------------------------------------------
% \role{title}{dates}{organization}{location}
\newcommand{\role}[4]{%
  \vspace{3pt}%
  \noindent\textbf{#1}\hfill\textcolor{muted}{#2}\par
  \noindent\textit{#3}\hfill\textcolor{muted}{\textit{#4}}\par
  \vspace{1pt}%
}
% \degree{degree}{field}{institution}{date}
\newcommand{\degree}[4]{%
  \noindent\textbf{#1}, #2\hfill\textcolor{muted}{#4}\par
  \noindent\textit{#3}\par\vspace{2pt}%
}
\newcommand{\skillrow}[2]{\noindent\textbf{#1:} #2\par\vspace{1.5pt}}
% \project{name}{description}{url}{short link label}
\newcommand{\project}[4]{\noindent\textbf{#1} --- #2 \href{#3}{\small\texttt{#4}}\par\vspace{1.5pt}}

\newlist{tight}{itemize}{1}
\setlist[tight]{leftmargin=13pt,topsep=1pt,itemsep=1pt,parsep=0pt,label=\textcolor{accent}{\small$\bullet$}}

\begin{document}

% ============================== HEADER ====================================
\begin{center}
  {\LARGE\bfseries Benjamin (Ben) Whewell}\\[3pt]
  {\color{accent}\bfseries AI/ML Scientist \textperiodcentered{} Machine Learning Engineering}\\[4pt]
  {\small Santa Fe, NM, USA \textperiodcentered{}
   \href{https://www.ben-whewell.com}{ben-whewell.com} \textperiodcentered{}
   \href{https://github.com/bwhewe-13}{github.com/bwhewe-13} \textperiodcentered{}
   \href{https://www.linkedin.com/in/ben-whewell-133935137/}{LinkedIn} \textperiodcentered{}
   \href{https://scholar.google.com/citations?user=e4PPMC0AAAAJ}{Google Scholar} \textperiodcentered{}
   \href{https://orcid.org/0000-0001-7826-5525}{ORCiD}}\\[3pt]
  {\footnotesize\itshape\color{muted} Public version --- phone number and email address withheld.
   Request the complete CV at \href{https://www.ben-whewell.com/resume.html}{ben-whewell.com/resume.html}.}
\end{center}

% ============================== SUMMARY ===================================
\section{Summary}
AI/ML scientist and engineer with a PhD in computational engineering and eight years of
research and production experience building machine learning systems for large-scale
scientific computing. Owns the full lifecycle: data ingestion and feature engineering
through training, evaluation, deployment into production codebases, and performance
monitoring with hands-on user adoption. Current work spans LLM and retrieval-augmented
generation (RAG) systems, reinforcement learning for sequential decision-making, and neural
surrogates that cut compute cost within a prescribed error tolerance. Strong Python and C++
engineering on Linux, GPU, and HPC infrastructure.

% ============================ SKILLS ======================================
\section{Technical Skills}
\skillrow{Languages}{Python, C++, SQL, R, Cython, Bash}
\skillrow{ML \& Deep Learning}{PyTorch, TensorFlow, Keras, Scikit-Learn, XGBoost, Optuna, OpenCV; neural
  operators (DeepONet), autoencoders, surrogate modeling, hyperparameter optimization}
\skillrow{LLMs \& NLP}{Hugging Face Transformers, LangChain, LangGraph, LlamaIndex, RAG pipelines,
  vector databases, embeddings, Ollama (local inference), Pydantic, prompt design and evaluation}
\skillrow{Reinforcement Learning}{Gymnasium, Stable-Baselines3, Ray RLlib, PettingZoo, multi-agent RL,
  custom environment design, reward shaping}
\skillrow{MLOps \& Deployment}{Docker, Kubernetes, MLflow, Weights \& Biases, TensorBoard, DVC, CI/CD,
  model monitoring, feature stores, FastAPI, Flask, Django}
\skillrow{Data \& Infrastructure}{NumPy, Pandas, PySpark, Hadoop, PostgreSQL, MySQL, SQLite, MongoDB,
  Snowflake, ETL; AWS, Azure, Git/GitHub, Linux/Unix, HPC and GPU clusters}

% ========================== EXPERIENCE ====================================
\section{Experience}

\role{AI/ML Scientist}{May 2026 -- Present}{Los Alamos National Laboratory}{Los Alamos, NM}
\begin{tight}
  \item Build and validate end-to-end ML pipelines --- ingestion, preprocessing, feature
    engineering, training, evaluation --- for large-scale multiphysics simulation workloads.
  \item Deploy models into production HPC software with performance monitoring, and drive
    adoption through direct support of scientific end users.
  \item Develop LLM- and RAG-based assistants over domain-specific technical corpora, covering
    retrieval, grounding, and response evaluation.
  \item Combine HPC and AI to accelerate computational science workflows and shorten the
    cycle from experiment to result.
\end{tight}

\role{Metropolis Postdoctoral Fellow}{February 2025 -- May 2026}{Los Alamos National Laboratory}{Los Alamos, NM}
\begin{tight}
  \item Trained reinforcement learning agents for sequential decision-making that reduces error
    in numerical simulations, including custom environment design and reward engineering.
  \item Applied ML-driven optimization to energy-group structures for nuclear data, shrinking
    data volume while holding simulation accuracy within tolerance.
  \item Prototyped LLM and RAG applications for technical document search and simulation
    workflows, and evaluated local versus hosted inference trade-offs.
  \item Awarded the competitive Nicholas C. Metropolis Postdoctoral Fellowship; conference paper
    award winner at the 2025 ANS Mathematics and Computational Methods conference.
\end{tight}

\role{Graduate Student Researcher}{September 2018 -- January 2025}{University of Notre Dame}{Notre Dame, IN}
\begin{tight}
  \item Designed and shipped an ML pipeline that replaces expensive linear-algebra kernels with
    learned surrogates (autoencoders, deep networks, deep operator networks), achieving data
    reduction in production simulations within a prescribed error tolerance.
  \item Built open-source scientific Python, Cython, and C++ libraries with published API
    documentation and regression tests, used by other researchers.
  \item Developed single-grid error estimation methodology for model-risk evaluation of
    high-dimensional solvers; published in a peer-reviewed verification and UQ journal.
\end{tight}

\role{Graduate Student Assistant}{June 2021 -- August 2021}{Lawrence Livermore National Laboratory}{Livermore, CA}
\begin{tight}
  \item Integrated parallel C++ production physics codes across Livermore's HPC systems.
  \item Coordinated across multiple developer teams and downstream users to reconcile
    code-base and user requirements.
\end{tight}

\role{Graduate Student Assistant}{May 2019 -- July 2019}{Los Alamos National Laboratory}{Los Alamos, NM}
\begin{tight}
  \item Built ML models to quantify experimental error and detect outliers in nuclear
    cross-section data using measurement features and covariances; results published in
    \textit{Nuclear Instruments and Methods in Physics Research A}.
\end{tight}

\role{Boilers and Turbines Operator}{June 2015 -- June 2018}{NRG Energy Power Station}{Middletown, CT}
\begin{tight}
  \item Operated and maintained safety-critical power-plant systems under strict procedures,
    with shift-handoff communication to supervision and plant personnel.
\end{tight}

% =========================== PROJECTS =====================================
\section{Selected Open-Source Projects}
\project{MyPDFLibrary}{RAG command-line assistant over PDF corpora: local LLM inference with Ollama,
  embedding-based vector retrieval, and grounded question answering.}{https://github.com/bwhewe-13/MyPDFLibrary}{bwhewe-13/MyPDFLibrary}
\project{WOLF}{Debugging and reliability contributor to Los Alamos National Laboratory's open-source
  agent framework (Python, FastAPI, Chroma): async coroutine leak, BM25 empty-corpus crash, PyMuPDF
  deprecation cleanup, and a missing knowledge-base PDF ingestion endpoint, root-caused and
  patched.}{https://github.com/lanl/wolf}{lanl/wolf}
\project{rl-lab}{Reinforcement learning library and tutorial series spanning tabular, deep,
  multi-agent, and distributed RL with PyTorch, Gymnasium, Stable-Baselines3, and Ray RLlib.}{https://github.com/bwhewe-13/rl-lab}{bwhewe-13/rl-lab}
\project{ants / discrete1}{Documented Python, Cython, and C++ simulation libraries with hooks for
  ML surrogate models and published API docs.}{https://github.com/bwhewe-13/ants}{bwhewe-13/ants}

% =========================== EDUCATION ====================================
\section{Education}
\degree{Doctor of Philosophy}{Mechanical and Aerospace Engineering}{University of Notre Dame}{January 2025}
\degree{Master of Science}{Mechanical and Aerospace Engineering}{University of Notre Dame}{May 2022}
\degree{Bachelor of Science}{Nuclear Engineering}{Thomas Edison State University}{December 2017}

% ========================= PUBLICATIONS ===================================
\section{Selected AI/ML Publications}
\begin{tight}
  \item \textbf{B. Whewell}, R. G. McClarren, N. Gibson, and A. Khatiwada (2026). Application of Deep
    Operator Networks for Data Reduction in Deterministic Neutron Calculations. \textit{Under review.}
  \item \textbf{B. Whewell}, N. Gibson, and A. Khatiwada (2026). Application of Reinforcement Learning for
    Multigroup Energy Grid Optimization for Neutron Transport Criticality Problems. \textit{Under review.}
  \item \textbf{B. Whewell} and R. G. McClarren (2022). Data Reduction in Deterministic Neutron Transport
    Calculations Using Machine Learning. \textit{Annals of Nuclear Energy}, 176, 109276.
    \href{https://doi.org/10.1016/j.anucene.2022.109276}{doi:10.1016/j.anucene.2022.109276}
  \item \textbf{B. Whewell}, M. Grosskopf, and D. Neudecker (2020). Evaluating $^{239}$Pu(n,f) cross sections
    via machine learning using experimental data, covariances, and measurement features.
    \textit{Nuclear Instruments and Methods in Physics Research A}, 978, 164305.
    \href{https://doi.org/10.1016/j.nima.2020.164305}{doi:10.1016/j.nima.2020.164305}
  \item \textbf{B. Whewell} and R. G. McClarren (2025). Single Grid Error Estimation for Neutron Transport
    Solvers. \textit{Journal of Verification, Validation and Uncertainty Quantification}, 10(2), 021001.
    \href{https://doi.org/10.1115/1.4069426}{doi:10.1115/1.4069426}
\end{tight}
\vspace{2pt}
{\small Full record --- 6 journal articles, 6 conference papers, 5 presentations, 2 technical reports:
\href{https://www.ben-whewell.com/research.html}{ben-whewell.com/research.html}.}

% ======================= AWARDS AND SERVICE ===============================
\section{Awards \& Professional Service}
\begin{tight}
  \item Nicholas C. Metropolis Postdoctoral Fellowship, Los Alamos National Laboratory.
  \item Conference Paper Award, 2025 ANS Mathematics and Computational Methods conference.
  \item Session Chair, Machine Learning and Artificial Intelligence I \& IV, 2025 International
    Conference on Mathematics and Computational Methods Applied to Nuclear Science and Engineering.
  \item Reviewer, \textit{Annals of Nuclear Energy} and ANS Mathematics and Computation Division
    (2025--present).
  \item Volunteer judge, Los Alamos National Laboratory 2025 Annual Student Symposium.
\end{tight}

\end{document}
